\documentclass[10pt,a4paper]{amsart}
\usepackage{amsmath,amssymb,amsthm,graphicx}
\usepackage{verbatim}
\newenvironment{reference}{\comment}{\endcomment}
%\newenvironment{reference}{}{}
\newenvironment{slogan}{\comment}{\endcomment}
\newenvironment{history}{\comment}{\endcomment}

% For commutative diagrams we use Xy-pic
\usepackage[all]{xy}

% We use 2cell for 2-commutative diagrams.
\xyoption{2cell}
% Theorem environments.
%
\theoremstyle{plain}
\newtheorem{theorem}[subsection]{Theorem}
\newtheorem{proposition}[subsection]{Proposition}
\newtheorem{lemma}[subsection]{Lemma}

\theoremstyle{definition}
\newtheorem{definition}[subsection]{Definition}
\newtheorem{example}[subsection]{Example}
\newtheorem{exercise}[subsection]{Exercise}
\newtheorem{situation}[subsection]{Situation}

\theoremstyle{remark}
\newtheorem{remark}[subsection]{Remark}
\newtheorem{remarks}[subsection]{Remarks}

\numberwithin{equation}{subsection}

% Macros
%
\def\lim{\mathop{\mathrm{lim}}\nolimits}
\def\colim{\mathop{\mathrm{colim}}\nolimits}
\def\Spec{\mathop{\mathrm{Spec}}}
\def\Hom{\mathop{\mathrm{Hom}}\nolimits}
\def\Ext{\mathop{\mathrm{Ext}}\nolimits}
\def\SheafHom{\mathop{\mathcal{H}\!\mathit{om}}\nolimits}
\def\SheafExt{\mathop{\mathcal{E}\!\mathit{xt}}\nolimits}
\def\Sch{\mathit{Sch}}
\def\Mor{\mathop{\mathrm{Mor}}\nolimits}
\def\Ob{\mathop{\mathrm{Ob}}\nolimits}
\def\Sh{\mathop{\mathit{Sh}}\nolimits}
\def\NL{\mathop{N\!L}\nolimits}
\def\CH{\mathop{\mathrm{CH}}\nolimits}
\def\proetale{{pro\text{-}\acute{e}tale}}
\def\etale{{\acute{e}tale}}
\def\QCoh{\mathit{QCoh}}
\def\Ker{\mathop{\mathrm{Ker}}}
\def\Im{\mathop{\mathrm{Im}}}
\def\Coker{\mathop{\mathrm{Coker}}}
\def\Coim{\mathop{\mathrm{Coim}}}

% Boxtimes
%
\DeclareMathSymbol{\boxtimes}{\mathbin}{AMSa}{"02}

%
% Macros for moduli stacks/spaces
%
\def\QCohstack{\mathcal{QC}\!\mathit{oh}}
\def\Cohstack{\mathcal{C}\!\mathit{oh}}
\def\Spacesstack{\mathcal{S}\!\mathit{paces}}
\def\Quotfunctor{\mathrm{Quot}}
\def\Hilbfunctor{\mathrm{Hilb}}
\def\Curvesstack{\mathcal{C}\!\mathit{urves}}
\def\Polarizedstack{\mathcal{P}\!\mathit{olarized}}
\def\Complexesstack{\mathcal{C}\!\mathit{omplexes}}
% \Pic is the operator that assigns to X its picard group, usage \Pic(X)
% \Picardstack_{X/B} denotes the Picard stack of X over B
% \Picardfunctor_{X/B} denotes the Picard functor of X over B
\def\Pic{\mathop{\mathrm{Pic}}\nolimits}
\def\Picardstack{\mathcal{P}\!\mathit{ic}}
\def\Picardfunctor{\mathrm{Pic}}
\def\Deformationcategory{\mathcal{D}\!\mathit{ef}}
\usepackage{cleveref}
\crefname{originalreference}{original reference}{original references}
\Crefname{originalreference}{Original reference}{Original references}
\makeatletter
\newlabel{situation-loc-chern}{{0GUJ}{source}{}{sourcepointer-1}{}}
\newlabel{situation-loc-chern@cref}{{[originalreference][0][]0GUJ}{[1][0][]source}{}{}{}}
\makeatother
\crefname{originaldefinition}{original definition}{original definitions}
\Crefname{originaldefinition}{Original definition}{Original definitions}
\makeatletter
\newlabel{definition-localized-chern}{{0FB5}{source}{}{sourcepointer-2}{}}
\newlabel{definition-localized-chern@cref}{{[originaldefinition][0][]0FB5}{[1][0][]source}{}{}{}}
\makeatother
\crefname{originallemma}{original lemma}{original lemmas}
\Crefname{originallemma}{Original lemma}{Original lemmas}
\makeatletter
\newlabel{lemma-independent-loc-chern}{{0FB2}{source}{}{sourcepointer-3}{}}
\newlabel{lemma-independent-loc-chern@cref}{{[originallemma][0][]0FB2}{[1][0][]source}{}{}{}}
\makeatother
\crefname{originalreference}{original reference}{original references}
\Crefname{originalreference}{Original reference}{Original references}
\makeatletter
\newlabel{remark-restriction-bivariant}{{0F9Z}{source}{}{sourcepointer-4}{}}
\newlabel{remark-restriction-bivariant@cref}{{[originalreference][0][]0F9Z}{[1][0][]source}{}{}{}}
\makeatother
\crefname{originallemma}{original lemma}{original lemmas}
\Crefname{originallemma}{Original lemma}{Original lemmas}
\makeatletter
\newlabel{flat-lemma-complex-and-divisor-blowup-pre}{{0GTY}{source}{}{sourcepointer-5}{}}
\newlabel{flat-lemma-complex-and-divisor-blowup-pre@cref}{{[originallemma][0][]0GTY}{[1][0][]source}{}{}{}}
\makeatother
\crefname{originallemma}{original lemma}{original lemmas}
\Crefname{originallemma}{Original lemma}{Original lemmas}
\makeatletter
\newlabel{flat-lemma-complex-and-divisor-blowup-base-change}{{0F8V}{source}{}{sourcepointer-6}{}}
\newlabel{flat-lemma-complex-and-divisor-blowup-base-change@cref}{{[originallemma][0][]0F8V}{[1][0][]source}{}{}{}}
\makeatother
\crefname{originallemma}{original lemma}{original lemmas}
\Crefname{originallemma}{Original lemma}{Original lemmas}
\makeatletter
\newlabel{lemma-localized-chern-pre}{{0F9K}{source}{}{sourcepointer-7}{}}
\newlabel{lemma-localized-chern-pre@cref}{{[originallemma][0][]0F9K}{[1][0][]source}{}{}{}}
\makeatother
\crefname{originallemma}{original lemma}{original lemmas}
\Crefname{originallemma}{Original lemma}{Original lemmas}
\makeatletter
\newlabel{lemma-localized-chern-pre-independent}{{0FAU}{source}{}{sourcepointer-8}{}}
\newlabel{lemma-localized-chern-pre-independent@cref}{{[originallemma][0][]0FAU}{[1][0][]source}{}{}{}}
\makeatother
\crefname{originallemma}{original lemma}{original lemmas}
\Crefname{originallemma}{Original lemma}{Original lemmas}
\makeatletter
\newlabel{lemma-base-change-localized-chern-pre}{{0GUI}{source}{}{sourcepointer-9}{}}
\newlabel{lemma-base-change-localized-chern-pre@cref}{{[originallemma][0][]0GUI}{[1][0][]source}{}{}{}}
\makeatother
\title{The localized Chern classes are compatible with base change of finite type}
\author{The Stacks Project Authors}
\date{stacks-project:chow.tex:local-snapshot-20261012; GNU FDL 1.2 or later}
\begin{document}
\maketitle
\noindent\textit{Editorial note.} All delivered mathematical text is the author's own, copied from the pinned Stacks Project chapter file chow.tex; no statement or proof step was rewritten, reordered or removed.\par\medskip
\noindent\textit{Editorial note.} Preamble provenance: the pinned chapter file opens with an input of the shared preamble and carries no preamble text of its own, so the project's single Stacks-wide preamble staged at sources/stacks-chow/preamble.tex is excerpted verbatim as two spans, lines 14-18 (Xy-pic setup) and lines 154-217 (the native theorem declarations and mathematical macros). Nothing was added, removed or redefined; the wrapper class and the amsmath/amssymb/amsthm/graphicx base packages are supplied by the assembler.\par\medskip
\noindent\textit{Editorial note.} The counted statement and proof contain no cross reference, no citation and no included file, so the delivered body is self contained.\par\medskip
\noindent\textit{Editorial note.} Canonical location of the counted proposition: Stacks tag 0FB3 (sources/stacks-chow/pinned-original-source.tex line 10824).\par\medskip
\noindent\textit{Allowed standard prerequisites (editorial).}\begin{itemize}
\item Construction of the localized Chern classes from the blowing up and the complex Q. Same chapter, Lemma labelled lemma-localized-chern-pre and Lemma labelled lemma-independent-loc-chern, cited at proof lines 10868-10873. Identifies the classes being compared with P'\_p(Q) and P'\_p(Q').
\item Base change of the blowing up complex and the description of Q as a derived twist. More on Flatness, Lemma labelled flat-lemma-complex-and-divisor-blowup-pre and Lemma labelled flat-lemma-complex-and-divisor-blowup-base-change, cited at proof lines 10843-10865. Provides the commutative square with Q' = (q o g)\textasciicircum{}*Q that realises the base change of the construction.
\item Base change of the localized class itself and independence of the chosen closed subscheme T'. Same chapter, Lemma labelled lemma-base-change-localized-chern-pre, cited at proof lines 10887-10893, and Lemma labelled lemma-localized-chern-pre-independent, cited at proof lines 10882-10886 and 10897-10900. Compares the restriction of P'\_p(Q) with the class built from q\textasciicircum{}*Q and shows P'\_p(Q') = P'\_p(q\textasciicircum{}*Q).
\item Restriction of bivariant classes. Same chapter, Remark labelled remark-restriction-bivariant, invoked in the statement. Fixes the meaning of the comparison asserted by the lemma.
\end{itemize}\medskip
\ifdefined\hypertarget\hypertarget{sourcepointer-1}{}\fi
\noindent\textit{Source reference.} Original reference 0GUJ: sources/stacks-chow/pinned-original-source.tex, source lines 10678--10678. Same-chapter pointer into the author environment situation at source line 10678; invoked at proof line None of the counted proof. It is retained only as a navigation pointer; every load-bearing step is inside the delivered span.\par\medskip
\ifdefined\hypertarget\hypertarget{sourcepointer-2}{}\fi
\noindent\textit{Source reference.} Original definition 0FB5: sources/stacks-chow/pinned-original-source.tex, source lines 10787--10787. Same-chapter pointer into the author environment definition at source line 10787; invoked at proof line None of the counted proof. It is retained only as a navigation pointer; every load-bearing step is inside the delivered span.\par\medskip
\ifdefined\hypertarget\hypertarget{sourcepointer-3}{}\fi
\noindent\textit{Source reference.} Original lemma 0FB2: sources/stacks-chow/pinned-original-source.tex, source lines 10701--10701. Same-chapter pointer into the author environment lemma at source line 10701; invoked at proof line None of the counted proof. It is retained only as a navigation pointer; every load-bearing step is inside the delivered span.\par\medskip
\ifdefined\hypertarget\hypertarget{sourcepointer-4}{}\fi
\noindent\textit{Source reference.} Original reference 0F9Z: sources/stacks-chow/pinned-original-source.tex, source lines 6184--6184. Same-chapter pointer into the author environment remark at source line 6184; invoked at proof line None of the counted proof. It is retained only as a navigation pointer; every load-bearing step is inside the delivered span.\par\medskip
\ifdefined\hypertarget\hypertarget{sourcepointer-5}{}\fi
\noindent\textit{Source reference.} Original lemma 0GTY: sources/stacks-chow/tags.txt, source lines 20983--20983. Cross-chapter standard prerequisite located in the project tag list staged in the local snapshot; invoked at proof line 10848 of the counted proof. It is retained only as a navigation pointer; every load-bearing step is inside the delivered span.\par\medskip
\ifdefined\hypertarget\hypertarget{sourcepointer-6}{}\fi
\noindent\textit{Source reference.} Original lemma 0F8V: sources/stacks-chow/tags.txt, source lines 19055--19055. Cross-chapter standard prerequisite located in the project tag list staged in the local snapshot; invoked at proof line 10861 of the counted proof. It is retained only as a navigation pointer; every load-bearing step is inside the delivered span.\par\medskip
\ifdefined\hypertarget\hypertarget{sourcepointer-7}{}\fi
\noindent\textit{Source reference.} Original lemma 0F9K: sources/stacks-chow/pinned-original-source.tex, source lines 10204--10204. Same-chapter pointer into the author environment lemma at source line 10204; invoked at proof line 10877 of the counted proof. It is retained only as a navigation pointer; every load-bearing step is inside the delivered span.\par\medskip
\ifdefined\hypertarget\hypertarget{sourcepointer-8}{}\fi
\noindent\textit{Source reference.} Original lemma 0FAU: sources/stacks-chow/pinned-original-source.tex, source lines 10301--10301. Same-chapter pointer into the author environment lemma at source line 10301; invoked at proof line 10886 of the counted proof. It is retained only as a navigation pointer; every load-bearing step is inside the delivered span.\par\medskip
\ifdefined\hypertarget\hypertarget{sourcepointer-9}{}\fi
\noindent\textit{Source reference.} Original lemma 0GUI: sources/stacks-chow/pinned-original-source.tex, source lines 10276--10276. Same-chapter pointer into the author environment lemma at source line 10276; invoked at proof line 10889 of the counted proof. It is retained only as a navigation pointer; every load-bearing step is inside the delivered span.\par\medskip
\begin{lemma}
\label{lemma-base-change-loc-chern}
In Situation \ref{situation-loc-chern}
let $f : X' \to X$ be a morphism of schemes which is locally of finite type.
Denote $E' = f^*E$ and $Z' = f^{-1}(Z)$. Then the bivariant class
of Definition \ref{definition-localized-chern}
$$
P_p(Z' \to X', E') \in A^p(Z' \to X'),
\quad\text{resp.}\quad
c_p(Z' \to X', E') \in A^p(Z' \to X')
$$
constructed as in Lemma \ref{lemma-independent-loc-chern}
using $X', Z', E'$ is the restriction
(Remark \ref{remark-restriction-bivariant}) of the
bivariant class $P_p(Z \to X, E) \in A^p(Z \to X)$,
resp.\ $c_p(Z \to X, E) \in A^p(Z \to X)$.
\end{lemma}

\begin{proof}
Denote $p : \mathbf{P}^1_X \to X$ and $p' : \mathbf{P}^1_{X'} \to X'$
the structure morphisms.
Recall that $b : W \to \mathbf{P}^1_X$ and $b' : W' \to \mathbf{P}^1_{X'}$
are the morphism constructed from the triples
$(\mathbf{P}^1_X, (\mathbf{P}^1_X)\infty, p^*E)$ and
$(\mathbf{P}^1_{X'}, (\mathbf{P}^1_{X'})\infty, (p')^*E')$
in More on Flatness, Lemma \ref{flat-lemma-complex-and-divisor-blowup-pre}.
Furthermore $Q = L\eta_{\mathcal{I}_\infty}p^*E$ and
$Q = L\eta_{\mathcal{I}'_\infty}(p')^*E'$
where
$\mathcal{I}_\infty \subset \mathcal{O}_W$ is the
ideal sheaf of $W_\infty$ and
$\mathcal{I}'_\infty \subset \mathcal{O}_{W'}$ is the
ideal sheaf of $W'_\infty$.
Next, $h : \mathbf{P}^1_{X'} \to \mathbf{P}^1_X$ is a morphism of
schemes such that the pullback of the effective Cartier divisor
$(\mathbf{P}^1_X)_\infty$ is the effective Cartier divisor
$(\mathbf{P}^1_{X'})_\infty$ and such that $h^*p^*E = (p')^*E'$.
By More on Flatness, Lemma
\ref{flat-lemma-complex-and-divisor-blowup-base-change}
we obtain a commutative diagram
$$
\xymatrix{
W' \ar[rd]_{b'} \ar[r]_-g &
\mathbf{P}^1_{X'} \times_{\mathbf{P}^1_X} W \ar[d]_r \ar[r]_-q &
W \ar[d]^b \\
&
\mathbf{P}^1_{X'} \ar[r] &
\mathbf{P}^1_X
}
$$
such that $W'$ is the ``strict transform'' of $\mathbf{P}^1_{X'}$
with respect to $b$ and such that $Q' = (q \circ g)^*Q$.
Now recall that $P_p(Z \to X, E) = P'_p(Q)$,
resp.\ $c_p(Z \to X, E) = c'_p(Q)$ where $P'_p(Q)$, resp.\ $c'_p(Q)$
are constructed in Lemma \ref{lemma-localized-chern-pre}
using $b, Q, T'$ where $T'$ is a closed subscheme $T' \subset W_\infty$
with the following two properties:
(a) $T'$ contains all points of $W_\infty$ lying over $X \setminus Z$,
and (b) $Q|_{T'}$ is zero, resp.\ isomorphic to a finite locally free
module of rank $< p$ placed in degree $0$.
In the construction of Lemma \ref{lemma-localized-chern-pre}
we chose a particular closed subscheme $T'$ with properties (a) and (b)
but the precise choice of $T'$ is immaterial, see
Lemma \ref{lemma-localized-chern-pre-independent}.

\medskip\noindent
Next, by Lemma \ref{lemma-base-change-localized-chern-pre}
the restriction of the bivariant class $P_p(Z \to X, E) = P'_p(Q)$,
resp.\ $c_p(Z \to X, E) = c_p(Q')$
to $X'$ corresponds to the class $P'_p(q^*Q)$, resp.\ $c'_p(q^*Q)$
constructed as in Lemma \ref{lemma-localized-chern-pre} using
$r : \mathbf{P}^1_{X'} \times_{\mathbf{P}^1_X} W \to \mathbf{P}^1_{X'}$,
the complex $q^*Q$, and the inverse image $q^{-1}(T')$.

\medskip\noindent
Now by the second statement of
Lemma \ref{lemma-localized-chern-pre-independent}
we have $P'_p(Q') = P'_p(q^*Q)$, resp.\ $c'_p(q^*Q) = c'_p(Q')$.
Since $P_p(Z' \to X', E') = P'_p(Q')$, resp.\ $c_p(Z' \to X', E') = c'_p(Q')$
we conclude that the lemma is true.
\end{proof}

\end{document}
